\chapter[Leçon 35: Nei Leit kenne léieren \& Alltagsausdréck]{Leçon 35: Nei Leit kenne léieren, Alltagsausdréck a Vokabelen}

Dës Lektioun baséiert op engem interaktive Cours a behandalt wichteg Ausdréck aus dem Alldag, wéi zum Beispill wéi een \textbf{nei Leit kenne léiert}, wéi ee sech gutt mat den \textbf{Nopere} versteet, an néideg Wierder wéi \textbf{de Fächer} (éventail / hand fan). Zousätzlech gi mir duerch d'Prononziatioun vun den Diphthongen (\textit{éi}, \textit{äi}, \textit{ei}) mat Lautschrëft (IPA), d'Konjugatioun vum Verb \textit{sech verstoen}, d'Verb \textit{verschampeléieren} an 10 Geleeënheeten a Froen iwwer sozial Kontakter a Fräizäit.

\textit{Cette leçon est basée sur un cours interactif et traite d'expressions essentielles du quotidien, telles que comment \textbf{faire la connaissance de nouvelles personnes}, bien s'entendre avec ses \textbf{voisins}, ainsi que du vocabulaire utile comme \textbf{de Fächer} (l'éventail). Nous abordons également la prononciation des diphtongues (\textit{éi}, \textit{äi}, \textit{ei}) avec la transcription phonétique (API), le verbe réfléchi \textit{sech verstoen}, le verbe \textit{verschampeléieren} et 10 questions pratiques sur les contacts sociaux et les loisirs.}

\textit{This lesson is based on an interactive class session and covers essential everyday expressions, such as how to \textbf{meet new people}, get along well with your \textbf{neighbors}, and useful vocabulary like \textbf{de Fächer} (the hand fan). Additionally, we review the pronunciation of diphthongs (\textit{éi}, \textit{äi}, \textit{ei}) with IPA phonetic transcriptions, the reflexive verb \textit{sech verstoen}, the verb \textit{verschampeléieren}, and 10 practical questions on social contacts and leisure activities.}

\vspace{1.5em}

\section[Vokabelen an Ausdréck / Vocabulary \& Expressions]{Vokabelen an Ausdréck aus dem Cours \\/ Vocabulaire et expressions du cours \\/ Key Vocabulary and Expressions from the Lesson}

\begin{itemize}
    \item \textbf{de Fächer} (m., Pl. \textit{d'Fächer} / \textit{d'Fächeren}):
    \begin{itemize}
        \item \textit{Bedeitung / Meaning:} L'éventail / Hand fan (objet pour se rafraîchir).
        \item \textit{Wichteg Notiz / Important Note:} \textbf{d'Fächer} ass och d'Méizuel (Plural) vum neutrum Substantiv \textbf{dat Fach} (n., Pl. \textit{d'Fächer} = la matière scolaire, le compartiment / school subject, compartment).
        \item \textit{Français :} L'éventail. (Note : \textit{d'Fächer} est aussi le pluriel de \textit{dat Fach} = la matière scolaire).
        \item \textit{English:} Hand fan. (Note: \textit{d'Fächer} is also the plural of \textit{dat Fach} = school subject/compartment).
        \item \textit{Beispill 1 (Fächer):} \textit{Wann et am Summer ganz waarm ass, benotzt si gär e Fächer.}
        $\rightarrow$ \textit{Français :} Quand il fait très chaud en été, elle aime utiliser un éventail.
        $\rightarrow$ \textit{English:} When it is very warm in summer, she likes using a hand fan.
        \item \textit{Beispill 2 (Fach):} \textit{Ech hunn zwee nei Fächer op der Schoul dëst Trimester.}
        $\rightarrow$ \textit{Français :} J'ai deux nouvelles matières à l'école ce trimestre.
        $\rightarrow$ \textit{English:} I have two new subjects at school this term.
    \end{itemize}

    \vspace{0.6em}

    \item \textbf{den Noper} / \textbf{d'Nopesch} / \textbf{d'Noperschaft}:
    \begin{itemize}
        \item \textbf{den Noper} (m., Pl. \textit{d'Noperen}): Le voisin / Male neighbor. (Opgepasst: wéinst dem Ufanksbuschtaf « N » bleift den « n » vum Artikel erhalen: \textbf{den} Noper).
        \item \textbf{d'Nopesch} (f., Pl. \textit{d'Nopeschen}): La voisine / Female neighbor.
        \item \textbf{d'Noperschaft} (f., Pl. \textit{d'Noperschaften}): Le voisinage, les voisins / Neighborhood, neighbors.
        \item \textit{Beispill:} \textit{Si verstinn sech ganz gutt mat hiren Noperen.}
        $\rightarrow$ \textit{Français :} Ils s'entendent très bien avec leurs voisins.
        $\rightarrow$ \textit{English:} They get along very well with their neighbors.
    \end{itemize}

    \vspace{0.6em}

    \item \textbf{verschampeléieren} (transitiivt Verb, verschampeléiert / verschampeléieren):
    \begin{itemize}
        \item \textit{Français :} Défigurer, enlaidir, abîmer l'aspect de.
        \item \textit{English:} To disfigure, to spoil the appearance of, to ruin.
        \item \textit{Synonymen:} entstellen, verschandelen.
        \item \textit{Aussprooch:} [fɐʃɑmpəlˈɜɪəʀən]
        \item \textit{Beispill 1:} \textit{Déi elle Stroosseluuchte verschampeléieren d'Landschaft.}
        $\rightarrow$ \textit{Français :} Ces lampadaires laids défigurent le paysage.
        $\rightarrow$ \textit{English:} Those ugly streetlights spoil the landscape.
        \item \textit{Beispill 2:} \textit{Zënter sengem Accident ass säi Gesiicht ganz verschampeléiert.}
        $\rightarrow$ \textit{Français :} Depuis son accident, son visage est tout défiguré.
        $\rightarrow$ \textit{English:} Since his accident, his face is completely disfigured.
    \end{itemize}

    \item \textbf{nei Leit kenne léieren} (léiert kennen / kenne geléiert):
    \begin{itemize}
        \item \textit{Français :} Faire la connaissance de nouvelles personnes, rencontrer du monde.
        \item \textit{English:} To meet new people, to make new acquaintances.
        \item \textit{Beispill:} \textit{Ech gi gär op oppen Eventer, well ech do einfach nei Leit kennen ze léiere sichen.}
        $\rightarrow$ \textit{Français :} J'aime aller à des événements ouverts pour rencontrer facilement de nouvelles personnes.
        $\rightarrow$ \textit{English:} I like going to open events to easily meet new people.
    \end{itemize}

    \vspace{0.6em}

    \item \textbf{beräicherend} (Adjektiv):
    \begin{itemize}
        \item \textit{Français :} Enrichissant, bénéfique.
        \item \textit{English:} Enriching, rewarding.
        \item \textit{Beispill:} \textit{Nei Leit kennen ze léieren ass meeschtens ganz beräicherend.}
        $\rightarrow$ \textit{Français :} Faire la connaissance de nouvelles personnes est généralement très enrichissant.
        $\rightarrow$ \textit{English:} Getting to know new people is usually very enriching.
    \end{itemize}

    \vspace{0.6em}

    \item \textbf{de Ventilator} (m., Pl. \textit{d'Ventilatoren}):
    \begin{itemize}
        \item \textit{Français :} Le ventilateur.
        \item \textit{English:} The ventilator, electric fan.
        \item \textit{Beispill:} \textit{Wann de Fächer net duergeet, maache mer de Ventilator un.}
        $\rightarrow$ \textit{Français :} Si l'éventail ne suffit pas, nous allumons le ventilateur.
        $\rightarrow$ \textit{English:} If the hand fan is not enough, we turn on the electric fan.
    \end{itemize}

    \vspace{0.6em}

    \item \textbf{krank geschriwwe sinn} / \textbf{e Certificat médical hunn}:
    \begin{itemize}
        \item \textit{Français :} Être en arrêt maladie / Avoir un certificat médical.
        \item \textit{English:} To be signed off sick / To have a medical certificate.
        \item \textit{Beispill:} \textit{Hien ass haut net op der Aarbecht, well hie vum Dokter krank geschriwwen ass.}
        $\rightarrow$ \textit{Français :} Il n'est pas au travail aujourd'hui parce qu'il est en arrêt maladie du médecin.
        $\rightarrow$ \textit{English:} He is not at work today because he was signed off sick by the doctor.
    \end{itemize}

    \vspace{0.6em}

    \item \textbf{d'Nohëllef} (f., Pl. \textit{d'Nohëllefen}) / \textbf{den Nohëllefsunterricht} (m.):
    \begin{itemize}
        \item \textit{Français :} Le soutien scolaire, les cours particuliers.
        \item \textit{English:} Tutoring, academic support, private lessons.
        \item \textit{Beispill:} \textit{Si gëtt de Kanner gären Nohëllef an de Sproochen.}
        $\rightarrow$ \textit{Français :} Elle donne volontiers du soutien scolaire aux enfants en langues.
        $\rightarrow$ \textit{English:} She enjoys giving children tutoring in languages.
    \end{itemize}
\end{itemize}

\vspace{1.5em}

\section[Diphthongen éi, äi, ei / Diphthongs éi, äi, ei]{Grammaire \& Prononziatioun: D'Diphthongen éi, äi an ei \\/ Prononciation et orthographe des diphtongues (avec IPA) \\/ Pronunciation and Spelling of Diphthongs (with IPA)}

\noindent E wichtegt Thema an der Lëtzebuerger Aussprooch a Schreifweis ass den Ënnerscheed tëscht den dräi Vokalkombinatiounen \textbf{éi}, \textbf{äi} an \textbf{ei}. Hei sinn d'Regelen an déi offiziell Lautschrëft (IPA):

\begin{enumerate}
    \item \textbf{éi} [əɪ] / [eːi̯] (zouenen, laangen « é-i » Diphthong):
    \begin{itemize}
        \item \textit{Français :} Prononcé comme un « é-i » fermé et allongé.
        \item \textit{English:} Pronounced as a closed long "e-i" diphthong sound.
        \item \textit{Wierder mat IPA:}
        \begin{itemize}
            \item \textbf{schéin} [ʃəɪn] (beau, joli / beautiful, nice)
            \item \textbf{séier} [ˈzeːi̯əʀ] (rapide, vite / fast, quick)
            \item \textbf{léieren} [ˈleːi̯əʀən] (apprendre, enseigner / to learn, to teach)
        \end{itemize}
    \end{itemize}

    \vspace{0.6em}

    \item \textbf{äi} [æːɪ] (oppenen « a-i » Sound):
    \begin{itemize}
        \item \textit{Français :} Prononcé comme un « a-i » ouvert.
        \item \textit{English:} Pronounced as an open "a-i" diphthong sound.
        \item \textit{Wierder mat IPA:}
        \begin{itemize}
            \item \textbf{Wäin} [væːɪn] (le vin / wine)
            \item \textbf{Schäin} [ʃæːɪn] (le certificat, la lumière / certificate, light)
            \item \textbf{bäi} [bæːɪ] (près de, chez / near, by)
        \end{itemize}
    \end{itemize}

    \vspace{0.6em}

    \item \textbf{ei} [ɑɪ] (klassichen Diphthong « ei »):
    \begin{itemize}
        \item \textit{Français :} Prononcé comme le diphtongue classique « ey/ai ».
        \item \textit{English:} Pronounced like the standard English "ey/ai" diphthong sound.
        \item \textit{Wierder mat IPA:}
        \begin{itemize}
            \item \textbf{Leit} [lɑɪt] (les gens, les personnes / people, folks)
            \item \textbf{nei} [nɑɪ] (neuf, nouveau / new, recent)
            \item \textbf{schreiwen} [ˈʃʀɑɪvən] (écrire / to write)
        \end{itemize}
    \end{itemize}
\end{enumerate}

\vspace{1.5em}

\section[Mündlech Ausdréck / Oral Expressions]{Mündlech Ausdréck: Kontakter an nei Leit kenne léieren \\/ Expressions orales : Contacts sociaux et loisirs \\/ Oral Expressions: Social Contacts and Meeting People}

\begin{enumerate}
    \item \textbf{Léiert Dir gär nei Leit kennen?}
    \begin{itemize}
        \item \textit{Lëtzebuergesch:} Jo, ech léiere ganz gär nei Leit kennen, well dat meeschtens ganz beräicherend ass a mäin Horizont erweidert.
        \item \textit{Français :} Oui, j'aime beaucoup faire la connaissance de nouvelles personnes, car c'est généralement très enrichissant et cela élargit mes horizons.
        \item \textit{English:} Yes, I really enjoy meeting new people because it is usually very enriching and broadens my horizons.
    \end{itemize}

    \item \textbf{Wou kënnt Dir nei Leit zu Lëtzebuerg kenne léieren?}
    \begin{itemize}
        \item \textit{Lëtzebuergesch:} Een kann nei Leit an de Vereiner, an de Sproochecoursen, bei oppenen Eventer oder beim Sport kenne léieren.
        \item \textit{Français :} On peut rencontrer de nouvelles personnes dans les clubs, les cours de langues, les événements ouverts ou le sport.
        \item \textit{English:} One can meet new people in clubs, language courses, open events, or sports.
    \end{itemize}

    \item \textbf{Wat maacht Dir, wann et am Summer ganz waarm ass?}
    \begin{itemize}
        \item \textit{Lëtzebuergesch:} Wann et ganz waarm ass, maachen ech d'Fënsteren zou, benotzen de Fächer oder de Ventilator an drénke vill Waasser.
        \item \textit{Français :} Quand il fait très chaud, je ferme les fenêtres, j'utilise un éventail ou un ventilateur et je bois beaucoup d'eau.
        \item \textit{English:} When it is very hot, I close the windows, use a hand fan or ventilator, and drink plenty of water.
    \end{itemize}

    \item \textbf{Versteet Dir Iech gutt mat Ären Noperen?}
    \begin{itemize}
        \item \textit{Lëtzebuergesch:} Jo, ech versti mech ganz gutt mat mengen Noperen, mir begréissen eis ëmmer an hëllefen eis wann et néideg ass.
        \item \textit{Français :} Oui, je m'entends très bien avec mes voisins, nous nous saluons toujours et nous nous entraidons si nécessaire.
        \item \textit{English:} Yes, I get along very well with my neighbors; we always greet each other and help out when needed.
    \end{itemize}

    \item \textbf{Sidd Dir schonn eng Kéier krank geschriwwe gewiescht?}
    \begin{itemize}
        \item \textit{Lëtzebuergesch:} Jo, wann ech eng Gripp hat, sinn ech vum Dokter krank geschriwwe ginn an hu mäi Certificat médical der Gemeng ginn.
        \item \textit{Français :} Oui, quand j'avais la grippe, le médecin m'a mis en arrêt maladie et j'ai transmis mon certificat médical.
        \item \textit{English:} Yes, when I had the flu, I was signed off sick by the doctor and submitted my medical certificate.
    \end{itemize}

    \item \textbf{Gitt Dir gär an de Sproochecours?}
    \begin{itemize}
        \item \textit{Lëtzebuergesch:} Jo, ech gi ganz gär an de Cours, well d'Atmosphär flott ass an een do zesumme Lëtzebuergesch übe kann.
        \item \textit{Français :} Oui, j'aime beaucoup aller au cours de langue parce que l'atmosphère est agréable et qu'on peut pratiquer le luxembourgeois ensemble.
        \item \textit{English:} Yes, I really enjoy going to language class because the atmosphere is nice and we can practice Luxembourgish together.
    \end{itemize}

    \item \textbf{Wat fir eng Aussprooch fënnt Dir schwéier am Lëtzebuergeschen?}
    \begin{itemize}
        \item \textit{Lëtzebuergesch:} Den Ënnerscheed tëscht den Diphthonge wéi « éi » [ʃəɪn] an « äi » [væːɪn] brauch e bëssen Übung, mä mat der Zäit klappt dat ëmmer besser.
        \item \textit{Français :} La différence entre les diphtongues comme « éi » et « äi » demande un peu de pratique, mais avec le temps ça va de mieux en mieux.
        \item \textit{English:} The difference between diphthongs like "éi" and "äi" takes a bit of practice, but over time it improves.
    \end{itemize}

    \item \textbf{Hutt Dir schonn eng Kéier Nohëllef ginn oder kritt?}
    \begin{itemize}
        \item \textit{Lëtzebuergesch:} Jo, ech hu Kanner bei den Hausaufgaben a Sprooche gehollef, dat ass eng ganz freedeg Aarbecht.
        \item \textit{Français :} Oui, j'ai aidé des enfants pour les devoirs et les langues, c'est un travail très joyeux.
        \item \textit{English:} Yes, I helped children with homework and languages; it is very joyful work.
    \end{itemize}

    \item \textbf{Wat maacht Dir gär an Ärer Fräizäit?}
    \begin{itemize}
        \item \textit{Lëtzebuergesch:} A menger Fräizäit ginn ech gär kuerz spadséieren, liese Bicher a kache mat menge Frënn a mengem Noper.
        \item \textit{Français :} Pendant mon temps libre, j'aime me promener, lire des livres et cuisiner avec mes amis et mon voisin.
        \item \textit{English:} In my free time, I like going for short walks, reading books, and cooking with my friends and neighbor.
    \end{itemize}

    \item \textbf{Wéi bleift ee mat senge Frënn am Kontakt?}
    \begin{itemize}
        \item \textit{Lëtzebuergesch:} Mir telefonéiere reegelméisseg, schreiwe Messagen an treffen eis fir e Kaffi oder en Iessen.
        \item \textit{Français :} Nous téléphonons régulièrement, nous écrivons des messages et nous nous retrouvons pour un café ou un repas.
        \item \textit{English:} We call regularly, write messages, and meet up for coffee or a meal.
    \end{itemize}
\end{enumerate}

\vspace{1.5em}

\section[Grammaire / Grammar]{Grammaire: Reflexiivt Verb « sech verstoen », Modalverben \& Eifeler Regel \\/ Grammaire : Verbe réfléchi « sech verstoen », verbes modaux et règle de l'Eifel \\/ Grammar: Reflexive Verb "sech verstoen", Modal Verbs, and Eifeler Rule}

\subsection*{1. Reflexiivt Verb: \textbf{sech verstoen mat} (S'entendre avec / To get along with)}
\begin{itemize}
    \item \textbf{Bedeitung / Meaning:} Dëst Verb gëtt benotzt fir interpersonal Bezéiungen ze beschreiwen (s'entendre bien/mal avec qn / to get along well/poorly with sb).
    \item \textit{Konjugatioun am Präsens (Present Tense):}
    \begin{itemize}
        \item \textit{ech versti mech} (je m'entends / I get along)
        \item \textit{du verstees dech} (tu t'entends / you get along)
        \item \textit{hien / hatt / si / ee \textbf{versteet} sech} (il/elle s'entend / he/she gets along)
        \item \textit{mir verstinn eis} (nous nous entendons / we get along)
        \item \textit{dir versteet iech} (vous vous entendez / you get along)
        \item \textit{si verstinn sech} (ils/elles s'entendent / they get along)
    \end{itemize}
    \item \textit{Beispiller:}
    \begin{itemize}
        \item \textit{Mir verstinn eis ganz gutt mat eisen Noperen.} (Nous nous entendons très bien avec nos voisins. / We get along very well with our neighbors.)
        \item \textit{Hie versteet sech net mat jiddwerengem.} (Il ne s'entend pas avec tout le monde. / He does not get along with everyone.)
    \end{itemize}
\end{itemize}

\vspace{0.8em}

\subsection*{2. Infinitiv: Modalverben (ouni « ze ») vs. Prepositionalen Infinitiv (mat « ze »)}
\begin{itemize}
    \item \textbf{Modalverben} (\textit{kënnen}, \textit{müssen}, \textit{däerfen}, \textit{wëllen}, \textit{sollen}):
    \begin{itemize}
        \item Huelen ëmmer e \textbf{rengen Infinitiv ouni « ze »} (sans « ze » / bare infinitive without "ze").
        \item \textit{Beispill:} \textit{Wou kënnt Dir nei Leit \textbf{kenne léieren}?} (Ouni « ze »!)
        $\rightarrow$ \textit{Français :} Où pouvez-vous rencontrer de nouvelles personnes ?
        $\rightarrow$ \textit{English:} Where can you meet new people?
        \item \textit{Eifeler Regel Notiz:} Virum Buschtaf « l » fält den « n » vu \textit{kennen} ewech $\rightarrow$ \textit{kenne léieren}.
    \end{itemize}

    \vspace{0.4em}

    \item \textbf{Prepositionalen Infinitiv mat « ze »} (Finalsaz mat \textit{fir ... ze} oder no Verbe wéi \textit{probéieren}):
    \begin{itemize}
        \item Huelen den Infinitiv mat der Partikel \textbf{« ze »}.
        \item \textit{Beispill:} \textit{Ech gi gär op oppen Eventer, fir nei Leit \textbf{kennen ze léieren}.}
        $\rightarrow$ \textit{Français :} Je vais volontiers à des événements ouverts pour faire la connaissance de nouvelles personnes.
        $\rightarrow$ \textit{English:} I like going to open events to get to know new people.
        \item \textit{Eifeler Regel Notiz:} Virum Buschtaf « z » (vun « ze ») bleift den « n » vu \textit{kennen} erhalen (Reegel: n, d, t, z, h) $\rightarrow$ \textit{kennen ze léieren}.
    \end{itemize}
\end{itemize}

\vspace{0.8em}

\subsection*{3. Eifeler Regel N-Check (Lektioun 35)}
\begin{itemize}
    \item \textbf{den Noper} (net \textit{de Noper}): De Buschtaf \textbf{N} gehéiert zu den Ausnamebuschtawen (\textit{n, d, t, z, h}), dofir bleift den « n » vum Artikel erhalen.
    \item \textbf{Hie benotzt}: De Buschtaf \textbf{B} ass keen N-Reegel Buschtaf, dofir fält den « n » ewech.
    \item \textbf{e Fächer}: Virun \textbf{F} fält den « n » vum Indefinit-Artikel ewech.
    \item \textbf{e ganz schéinen Dag}: Viru \textbf{G} fält den « n » ewech.
    \item \textbf{Si hu sech}: Virum Buschtaf « s » (vun \textit{sech}) fält den « n » vun \textit{hunn} ewech.
    \item \textbf{Si si och}: Virum Buschtaf « s » (vun \textit{si}) fält den « n » vu \textit{sinn} ewech.
    \item \textbf{mat menge Frënn}: Virun \textbf{F} fält den « n » ewech.
\end{itemize}

\vspace{1.5em}

\section[Übungen / Exercises]{Übungen / Exercices / Exercises}

\subsection*{Übung 1: Wierder asetzen (Insert Vocabulary)}
\textit{Français : Complétez les phrases avec les mots suivants : \textbf{Fächer}, \textbf{kennen ze léieren}, \textbf{Nohëllef}, \textbf{Noperen}, \textbf{schéin}}\\
\textit{English: Fill in the blanks with the following words: \textbf{Fächer}, \textbf{kennen ze léieren}, \textbf{Nohëllef}, \textbf{Noperen}, \textbf{schéin}}

\begin{enumerate}
    \item Ech gi gär op oppen Eventer, fir nei Leit \rule{3.5cm}{0.4pt}.
    \item Wann et am Summer ganz waarm ass, benotzt si gär e \rule{3.5cm}{0.4pt}.
    \item Si ginn de Schüler gär \rule{3.5cm}{0.4pt} an de Sproochen.
    \item Dat ass e ganz \rule{3.5cm}{0.4pt} Dag haut fir e Spazéiergang.
    \item Ech versti mech ganz gutt mat mengen \rule{3.5cm}{0.4pt}.
\end{enumerate}

\subsection*{Übung 2: Iwwersetzung (Translation into Luxembourgish)}
\textit{Français : Traduisez les phrases suivantes en luxembourgeois.}\\
\textit{English: Translate the following sentences into Luxembourgish.}

\begin{enumerate}
    \item \textit{Français :} Je m'entends très bien avec mes amis.\\
    \textit{English:} I get along very well with my friends.\\
    $\rightarrow$ \rule{8cm}{0.4pt}

    \vspace{0.6em}

    \item \textit{Français :} Il utilise un éventail quand il fait chaud.\\
    \textit{English:} He uses a hand fan when it is warm.\\
    $\rightarrow$ \rule{8cm}{0.4pt}

    \vspace{0.6em}

    \item \textit{Français :} Elle donne du soutien scolaire aux enfants.\\
    \textit{English:} She gives tutoring to the children.\\
    $\rightarrow$ \rule{8cm}{0.4pt}
\end{enumerate}

\vspace{1.5em}

\subsection*{Léisungen (Solutions)}

\subsubsection*{Léisung fir Übung 1}
\begin{enumerate}
    \item Ech gi gär op oppen Eventer, fir nei Leit \textbf{kennen ze léieren}.
    $\rightarrow$ \textit{Français :} J'aime aller à des événements ouverts pour faire la connaissance de nouvelles personnes.
    $\rightarrow$ \textit{English:} I like going to open events to meet new people.

    \item Wann et am Summer ganz waarm ass, benotzt si gär e \textbf{Fächer}.
    $\rightarrow$ \textit{Français :} Quand il fait très chaud en été, elle aime utiliser un éventail.
    $\rightarrow$ \textit{English:} When it is very warm in summer, she likes using a hand fan.

    \item Si ginn de Schüler gär \textbf{Nohëllef} an de Sproochen.
    $\rightarrow$ \textit{Français :} Ils/Elles donnent volontiers du soutien scolaire aux élèves en langues.
    $\rightarrow$ \textit{English:} They enjoy giving students tutoring in languages.

    \item Dat ass e ganz \textbf{schéinen} Dag haut fir e Spazéiergang.
    $\rightarrow$ \textit{Français :} C'est une très belle journée aujourd'hui pour une promenade.
    $\rightarrow$ \textit{English:} It is a very beautiful day today for a walk.

    \item Ech versti mech ganz gutt mat mengen \textbf{Noperen}.
    $\rightarrow$ \textit{Français :} Je m'entends très bien avec mes voisins.
    $\rightarrow$ \textit{English:} I get along very well with my neighbors.
\end{enumerate}

\subsubsection*{Léisung fir Übung 2}
\begin{enumerate}
    \item \textit{Français :} Je m'entends très bien avec mes amis.
    $\rightarrow$ \textbf{Ech versti mech ganz gutt mat menge Frënn.}

    \item \textit{Français :} Il utilise un éventail quand il fait chaud.
    $\rightarrow$ \textbf{Hie benotzt e Fächer wann et waarm ass.}

    \item \textit{Français :} Elle donne du soutien scolaire aux enfants.
    $\rightarrow$ \textbf{Si gëtt de Kanner Nohëllef.}
\end{enumerate}

\vspace{2em}

\section[Vokabelen mat Biller / Picture Dictionary]{Vokabelen mat Biller / Dictionnaire illustré / Illustrated Dictionary}

\begin{noindent}
\begin{minipage}[t]{0.48\textwidth}
    \centering
    \textbf{\Large E Fächer}
    \vspace{0.4em}
    \framebox{\parbox{0.88\linewidth}{\centering
            \vspace{0.3cm}
            \includegraphics[width=0.85\linewidth]{vokab_300dpi/facher.png} \\
            \vspace{0.3cm}
    }}
    \vspace{0.4em}
    \textbf{\Large De Fächer}\\
    \small\textit{(L'éventail / The hand fan)}
\end{minipage}
\hfill
\begin{minipage}[t]{0.48\textwidth}
    \centering
    \textbf{\Large Eng Tass}
    \vspace{0.4em}
    \framebox{\parbox{0.88\linewidth}{\centering
            \vspace{0.3cm}
            \includegraphics[width=0.85\linewidth]{vokab_300dpi/Tass.png} \\
            \vspace{0.3cm}
    }}
    \vspace{0.4em}
    \textbf{\Large D'Tass}\\
    \small\textit{(La tasse / The cup)}
\end{minipage}
\end{noindent}

\vspace{2em}

\begin{noindent}
\begin{minipage}[t]{0.48\textwidth}
    \centering
    \textbf{\Large Eng Famill}
    \vspace{0.4em}
    \framebox{\parbox{0.88\linewidth}{\centering
            \vspace{0.3cm}
            \includegraphics[width=0.85\linewidth]{vokab_300dpi/famill.png} \\
            \vspace{0.3cm}
    }}
    \vspace{0.4em}
    \textbf{\Large D'Famill}\\
    \small\textit{(La famille / The family)}
\end{minipage}
\hfill
\begin{minipage}[t]{0.48\textwidth}
    \centering
    \textbf{\Large En Noper}
    \vspace{0.4em}
    \framebox{\parbox{0.88\linewidth}{\centering
            \vspace{0.3cm}
            \includegraphics[width=0.85\linewidth]{vokab_300dpi/noper.png} \\
            \vspace{0.3cm}
    }}
    \vspace{0.4em}
    \textbf{\Large Den Noper}\\
    \small\textit{(Le voisin / The neighbor)}
\end{minipage}
\end{noindent}
